% TOP_PROBLEM: {"id": 3288, "title": "Weak saturation of the cube in the clique", "category_id": 3, "category": {"id": 3, "name": "graph_theory", "display_name": "Graph Theory", "description": "Problems involving graphs, networks, and their properties.", "slug": "graph-theory", "order_index": 3, "created_at": "2026-07-31T15:26:25.671Z"}, "status": "open", "problem_number": "OPG-60013", "set_id": 12}
% FIRST_POSTED: 2026-10-02
% ATTEMPT_STATUS: unresolved
% UnsolvedMath ID 3288; source catalogue code OPG-60013
% !TeX program = lualatex
% GPT-6 Astra Ultra research attempt, 2026-10-02; independently checked partial work.
% Completion estimate: no numerical percentage assigned; see final status and literature audit.
\documentclass[11pt,a4paper]{article}
\usepackage[margin=25mm]{geometry}
\usepackage{fontspec}
\setmainfont{lmroman10-regular.otf}[BoldFont=lmroman10-bold.otf,
 ItalicFont=lmroman10-italic.otf,BoldItalicFont=lmroman10-bolditalic.otf]
\usepackage{amsmath,amssymb,mathtools}
\usepackage{unicode-math}
\setmathfont{latinmodern-math.otf}
\AtBeginDocument{\let\setminus\smallsetminus}
\usepackage{xurl,hyperref}
\hypersetup{colorlinks=true,urlcolor=blue}
\setcounter{secnumdepth}{0}
\setlength{\parindent}{0pt}
\setlength{\parskip}{5pt}
\setlength{\emergencystretch}{3em}
\begin{document}
\section*{Weak saturation of the cube in the clique}\label{3288}
{\small UnsolvedMath ID 3288 · OPG-60013 · Graph Theory · OpenGarden}
\subsection{Definitions and mathematical statement}
Let $K_n$ be the complete simple graph on the labelled vertex set
$\{1,\ldots,n\}$, where $n\ge1$. The three-dimensional cube graph $Q_3$
has vertex set $\{0,1\}^3$, with two binary strings adjacent when they
differ in exactly one coordinate. It has eight vertices and twelve
edges.

A spanning subgraph $F\subseteq K_n$ is weakly $Q_3$-saturated in $K_n$
if the missing edges can be ordered $e_1,\ldots,e_m$ so that, for each
$i$, the graph
\[
                   F_i=F+e_1+\cdots+e_i
\]
contains a copy of $Q_3$ using $e_i$ and having all its other edges in
$F_{i-1}$. Thus each insertion creates a new cube, with the newly
inserted edge belonging to that cube. The copy is not required to be
induced, and already present edges not belonging to it are harmless.
The order of insertions may be chosen to suit $F$.

Define
\[
 \operatorname{wsat}(K_n,Q_3)=
       \min\{|E(F)|:F\text{ is weakly }Q_3\text{-saturated in }K_n\}.
\]
The problem is to determine this integer exactly as a function of $n$,
by matching a construction with a universal lower bound. An initial
cube-free condition is not imposed on $F$. For $n<8$ no cube can appear,
so no missing edge can be inserted; consequently the value is
$\binom n2$ in that range. The substantive target is the exact value
for $n\ge8$, rather than a weak-saturation process whose host is itself
a larger cube.
\subsection{Short English statement}
What is the fewest initial edges in a complete graph from which all missing edges can be added, each addition completing a cube graph?
\subsection{Research attempt}
\paragraph{Scope and process.}
GPT-6 Astra Ultra, 2 October 2026. The canonical formulation above is retained. This attempt uses a primary-source literature audit, independent restricted proofs, and adversarial checks. Reproved results are not claimed as new. Numerical tests below are reproducible in \texttt{scripts/check\_attempt\_batch03\_extremal\_2026\_10\_02.py}.

\paragraph{Literature and scope.}
The PCMI 2025 presentation [S3] reports $\operatorname{wsat}(K_8,Q_3)=15$, $\operatorname{wsat}(K_9,Q_3)=17$, and the upper bound $2n-1$, with equality conjectured in general. We independently reproduce the eight-vertex value and give a complete extension proof. Initial graphs are not required to be cube-free. For $n<8$, no cube can appear at all, so the only weakly saturated graph is the complete host itself.

\paragraph{A finite exact certificate at eight vertices.}
Label vertices $0,\ldots,7$ by their three-bit binary expansions. Let $Q$ have edges joining labels whose exclusive-or has exactly one set bit. Put
\[
 F=(E(Q)\setminus\{01\})\cup\{03,05,06,12\}.
\]
This seed has fifteen edges. A legal sequence inserting the remaining thirteen edges is
\[
 24,27,16,34,07,17,35,36,25,14,01,47,56.
\]
For each insertion the script supplies a labelled copy of $Q_3$ whose other eleven edges are already present. It generates copies directly by all permutations of the canonical cube and checks the indicated condition by integer edge masks. Thus the certificate does not depend on an isomorphism library or a heuristic search.

For the lower bound, there are $8!/48=840$ distinct cube edge sets in $K_8$; the script also obtains this count without assuming the automorphism count. Every successful seed with at most fourteen edges has a first inserted edge. Its witnessing cube, minus that edge, consists of eleven seed edges. Relabelling the cube and using its edge transitivity sends these eleven edges to $E(Q)\setminus\{01\}$. Therefore it suffices to inspect this fixed eleven-edge set augmented by at most three of the other seventeen host edges. There are exactly
\[
 \sum_{j=0}^{3}\binom{17}{j}=834
\]
such seeds. The exhaustive computation finds that none closes to $K_8$.

Here closure means repeatedly adding any missing edge that completes a cube. This exhaustion is valid independently of the order of choices: once an edge is eligible, later insertions cannot make it ineligible unless it has already been added. Induction on a proposed legal sequence shows that every edge in any such sequence belongs to the final greedy closure. A proper fixed point thus certifies failure of every possible ordering, not merely the tested ordering. Together with the displayed seed this proves the finite value fifteen, subject to the transparent exhaustive certificate.

\paragraph{Extending the seed to every larger host.}
For $n\ge8$, take the seed $F$ on a fixed eight-vertex core. Join every other vertex to two fixed distinct core vertices $a,b$, and include no other edges. There are $15+2(n-8)=2n-1$ initial edges. First saturate the core with the certificate above.

For an outside vertex $v$ and a core vertex $c\notin\{a,b\}$, insert $vc$. Map one vertex of a cube to $v$, its three neighbours to $a,b,c$, and its remaining four vertices to four other distinct core vertices. All needed edges except $vc$ are present: the only cube edges incident with $v$ are $va,vb,vc$, and every edge entirely in the core exists. This joins every outside vertex to the whole core. Finally insert any edge between outside vertices $v,w$ by mapping adjacent cube vertices to $v,w$ and its other six vertices to distinct core vertices. All remaining cube edges lie in the completed core or between an outside vertex and the core. Thus every host edge is eventually inserted.

\paragraph{Verification and limits.}
Every vertex of a weakly saturated noncomplete host seed must initially have degree at least two: its first new incident edge needs the other two incident cube edges already present. This gives only $|E(F)|\ge n$, far below the constructed $2n-1$. The finite normalization checks all small seeds, including those already containing other cubes; it cannot be extrapolated from eight vertices to arbitrary $n$. The companion script verifies all 834 closures and every edge of the fifteen-edge certificate.

\paragraph{Final status.}
The exact eight-vertex case and the general $2n-1$ construction are checked. A matching lower bound for arbitrary $n\ge8$, or a better construction, is the first remaining gap. The catalogue determination is \textbf{UNRESOLVED} by this attempt.

\subsection{Sources}
\begin{itemize}
\item[S1] ULAM.AI and the UnsolvedMath Contributors, supplied problems.json, record 3288 (OPG-60013), OpenGarden collection. Catalogue identity and source transcription; CC BY 4.0.
\url{https://huggingface.co/datasets/ulamai/UnsolvedMath}.
\item[S2] Open Problem Garden, Weak saturation of the cube in the clique. Original problem and discussion; the mathematical formulation below is expanded from the supplied source record.
\url{http://www.openproblemgarden.org/op/weak_saturation_of_the_cube_in_the_clique}.
\item[S3] T. Trakulthongchai, D. Xu and K. Zhu, Weak saturation number of the 3-cube in the complete graph, PCMI Summer 2025 presentation. The finite values and upper bound are prior reported results, independently checked in part here.
\url{https://www.ias.edu/sites/default/files/Trakulthongchai_Xu_Zhu.pdf}.
\end{itemize}
\end{document}
